% GENERATED FILE. Edit canonical lessons, then run npm run generate:books.
% canonical-source-hash: fnv1a64:39e3eab20374df0d
% canonical-lessons: ZH-C372-fanzheng, ZH-C372-daodi, ZH-C372-tongshi, ZH-C372-yongyi, ZH-C372-youdian

\chapter{Anyway}
\label{ch:zh-anyway}
\hlchaptermodality{372}

\begin{chapteropening}
\textbf{By the end of this chapter:} \emph{I can say anyway, ask what on earth someone means, and name someone's strong points.}

Complete the last lesson of chapter 372: I can say anyway, ask what on earth someone means, and name someone's strong points.
\end{chapteropening}

\section[fǎnzhèng]{fǎnzhèng — anyway; in any case}

\label{lesson:ZH-C372-fanzheng}



\begin{quote}
Before the new one: say the Mandarin for to sum up, then the Mandarin for a principle you live by.
\end{quote}

\subsection*{You'll want to know: fǎnzhèng}
\emph{fǎnzhèng} — \textquotedblleft{}anyway; in any case\textquotedblright{}.

\begin{quote}
\emph{nǐ zuìhǎo bié zháojí fǎnzhèng hái yǒu shíjiān} — Try not to worry; there's still time anyway.
\end{quote}

Its characters come later; for now it is a word for the ear.

\begin{cousinweb}
\emph{fǎnzhèng} is built from \emph{fǎn} and \zh{正} 'proper', both inherited from Middle Chinese, \emph{pjonX} and \emph{tsyengH} (Guangyun, in Baxter's transcription), which glosses the first 'to turn over' and the second 'proper'. How the pair came to mean anyway is not settled.
\end{cousinweb}

\begin{grammarlens}[title={What you've built}]
One more word for giving reasons.
\end{grammarlens}

\subsection*{Your turn}
\begin{itemize}
\raggedright
  \item \emph{Say it:} \emph{fǎnzhèng}
  \item \emph{Say it:} \emph{fǎnzhèng}, once more
  \item \emph{Say it:} \emph{nǐ zuìhǎo bié zháojí fǎnzhèng hái yǒu shíjiān} — Try not to worry; there's still time anyway.
  \item \emph{Recall:} say the Mandarin for to sum up, then the Mandarin for a principle you live by, then say \emph{fǎnzhèng} again
\end{itemize}

\begin{tcolorbox}[breakable,colback=teal!4,colframe=teal!35!black,title={Before you move on}]
What does fǎnzhèng mean? (\textquotedblleft{}Anyway; in any case\textquotedblright{}.) Say it once more.
\end{tcolorbox}
\section[dàodǐ]{dàodǐ — on earth (in a question)}

\label{lesson:ZH-C372-daodi}



\begin{quote}
Before the new one: say the Mandarin for a principle you live by, then the Mandarin for anyway.
\end{quote}

\subsection*{You'll want to know: dàodǐ}
\emph{dàodǐ} — \textquotedblleft{}on earth (in a question)\textquotedblright{}.

\begin{quote}
\emph{nǐ dàodǐ xiǎng shuō shénme} — What on earth are you trying to say?
\end{quote}

Its characters come later; for now it is a word for the ear.

\begin{cousinweb}
\emph{dàodǐ} is built from \zh{到} 'to arrive' and \emph{dǐ}, both inherited from Middle Chinese, \emph{tawH} and \emph{tejX} (Guangyun, in Baxter's transcription), which glosses the first 'to arrive' and the second 'below'.
\end{cousinweb}

\begin{grammarlens}[title={What you've built}]
One more word for giving reasons.
\end{grammarlens}

\subsection*{Your turn}
\begin{itemize}
\raggedright
  \item \emph{Say it:} \emph{dàodǐ}
  \item \emph{Say it:} \emph{dàodǐ}, once more
  \item \emph{Say it:} \emph{nǐ dàodǐ xiǎng shuō shénme} — What on earth are you trying to say?
  \item \emph{Recall:} say the Mandarin for a principle you live by, then the Mandarin for anyway, then say \emph{dàodǐ} again
\end{itemize}

\begin{tcolorbox}[breakable,colback=teal!4,colframe=teal!35!black,title={Before you move on}]
What does dàodǐ mean? (\textquotedblleft{}On earth (in a question)\textquotedblright{}.) Say it once more.
\end{tcolorbox}
\section[\texorpdfstring{\zh{同时} (tóngshí)}{同时 (tóngshí)}]{\zh{同时} (tóngshí) — simultaneously; at the same moment}

\label{lesson:ZH-C372-tongshi}



\begin{quote}
Before the new one: say the Mandarin for anyway, then the Mandarin for on earth (in a question).
\end{quote}

\subsection*{You'll want to know: \zh{同时}}
\textbf{\zh{同时}} — \emph{tóngshí} — \textquotedblleft{}simultaneously; at the same moment\textquotedblright{}.

\begin{quote}
\textbf{\zh{他们两个人同时到了}} \emph{tāmen liǎng ge rén tóngshí dào le} — The two of them arrived at the same time.
\end{quote}

\begin{cousinweb}
\textbf{\zh{同时}} (\emph{tóngshí}) is built from \zh{同} 'the same' and \zh{时} 'a time', both inherited from Middle Chinese, \emph{duwng} and \emph{dzyi} (Guangyun, in Baxter's transcription), which glosses the first 'equal' and the second 'a time of day'.
\end{cousinweb}

\begin{grammarlens}[title={What you've built}]
One more word for giving reasons.
\end{grammarlens}

\subsection*{Your turn}
\begin{itemize}
\raggedright
  \item \emph{Say it:} \emph{tóngshí}
  \item \emph{Say it:} \emph{tóngshí}, once more
  \item \emph{Say it:} \emph{tāmen liǎng ge rén tóngshí dào le} — The two of them arrived at the same time.
  \item \emph{Recall:} say the Mandarin for anyway, then the Mandarin for on earth (in a question), then say \emph{tóngshí} again
\end{itemize}

\begin{tcolorbox}[breakable,colback=teal!4,colframe=teal!35!black,title={Before you move on}]
What does \zh{同时} mean? (\textquotedblleft{}Simultaneously; at the same moment\textquotedblright{}.) Say it once more.
\end{tcolorbox}
\section[\texorpdfstring{\zh{用意} (yòngyì)}{用意 (yòngyì)}]{\zh{用意} (yòngyì) — intention; the purpose behind it}

\label{lesson:ZH-C372-yongyi}



\begin{quote}
Before the new one: say the Mandarin for on earth (in a question), then the Mandarin for simultaneously.
\end{quote}

\subsection*{You'll want to know: \zh{用意}}
\textbf{\zh{用意}} — \emph{yòngyì} — \textquotedblleft{}intention; the purpose behind it\textquotedblright{}.

\begin{quote}
\textbf{\zh{虽然他没有说但是我明白他的用意}} \emph{suīrán tā méiyǒu shuō dànshì wǒ míngbai tā de yòngyì} — He didn't say so, but I understand what he meant by it.
\end{quote}

\begin{cousinweb}
\textbf{\zh{用意}} (\emph{yòngyì}) is built from \zh{用} 'to use' and \zh{意} 'intent', both inherited from Middle Chinese, \emph{yowngH} and \emph{'iH} (Guangyun, in Baxter's transcription), which glosses the first 'to use' and the second 'will'.
\end{cousinweb}

\begin{grammarlens}[title={What you've built}]
One more word for giving reasons.
\end{grammarlens}

\subsection*{Your turn}
\begin{itemize}
\raggedright
  \item \emph{Say it:} \emph{yòngyì}
  \item \emph{Say it:} \emph{yòngyì}, once more
  \item \emph{Say it:} \emph{suīrán tā méiyǒu shuō dànshì wǒ míngbai tā de yòngyì} — He didn't say so, but I understand what he meant by it.
  \item \emph{Recall:} say the Mandarin for on earth (in a question), then the Mandarin for simultaneously, then say \emph{yòngyì} again
\end{itemize}

\begin{tcolorbox}[breakable,colback=teal!4,colframe=teal!35!black,title={Before you move on}]
What does \zh{用意} mean? (\textquotedblleft{}Intention; the purpose behind it\textquotedblright{}.) Say it once more.
\end{tcolorbox}
\section[yōudiǎn]{yōudiǎn — strong point; merit}

\label{lesson:ZH-C372-youdian}



\begin{quote}
Before the new one: say the Mandarin for simultaneously, then the Mandarin for the purpose behind it.
\end{quote}

\subsection*{You'll want to know: yōudiǎn}
\emph{yōudiǎn} — \textquotedblleft{}strong point; merit\textquotedblright{}.

\begin{quote}
\emph{wǒ juéde zhè shì tā zuì dà de yōudiǎn} — I think that's his greatest strength.
\end{quote}

Its characters come later; for now it is a word for the ear.

\begin{cousinweb}
\emph{yōudiǎn} is built from \emph{yōu} and \zh{点} 'a point', both inherited from Middle Chinese, \emph{'juw} and \emph{temX} (Guangyun, in Baxter's transcription), which glosses the first 'abundant' and the second 'a dot'.
\end{cousinweb}

\begin{grammarlens}[title={What you've built}]
One more word for giving reasons.
\end{grammarlens}

\subsection*{Your turn}
\begin{itemize}
\raggedright
  \item \emph{Say it:} \emph{yōudiǎn}
  \item \emph{Say it:} \emph{yōudiǎn}, once more
  \item \emph{Say it:} \emph{wǒ juéde zhè shì tā zuì dà de yōudiǎn} — I think that's his greatest strength.
  \item \emph{Recall:} say the Mandarin for simultaneously, then the Mandarin for the purpose behind it, then say \emph{yōudiǎn} again
\end{itemize}

\begin{tcolorbox}[breakable,colback=teal!4,colframe=teal!35!black,title={Before you move on}]
What does yōudiǎn mean? (\textquotedblleft{}Strong point; merit\textquotedblright{}.) Say it once more.
\end{tcolorbox}
